\QHeading{第1题　口算。}
\par\noindent


\begin{oralgrid}
\OralItem{1}{25\times6}{150}
\OralItem{2}{180\div6}{30}
\OralItem{3}{30\times40}{1200}
\OralItem{4}{560\div70}{8}
\OralItem{5}{450\div9}{50}
\OralItem{6}{15\times6}{90}
\OralItem{7}{720\div80}{9}
\OralItem{8}{20\times50}{1000}
\OralItem{9}{360\div40}{9}
\OralItem{10}{18\times5}{90}
\OralItem{11}{60\times13}{780}
\OralItem{12}{480\div6}{80}\end{oralgrid}
\par

\QHeading{口算思路补充（对应上表题号）}

\Answer{上面列出得数，下面按小题给出不同计算思路。计数单位法适用于小数，整数题可比较凑整、拆分或商不变等方法。}
\MethodTitle{第1小题：$25\times 6$}

\MethodBlock{口算思路一：按数位拆分计算}{把一个因数拆成整十、整百数与余下部分之和。}{\begin{aligned}[t]&\phantom{={}}25\times 6\\
&=6\times (20+5)\\
&=6\times 20+6\times 5\\
&=120+30\\
&=150\end{aligned}}
\MethodTitle{第2小题：$180\div 6$}

\MethodBlock{口算思路一：利用商不变的性质}{被除数和除数同时乘同一个非零数，商不变。}{\begin{aligned}[t]&\phantom{={}}180\div 6\\
&=(180\times 5)\div (6\times 5)\\
&=900\div 30\\
&=30\end{aligned}}

\MethodBlock{口算思路二：拆分除数，连续相除}{}{\begin{aligned}[t]&\phantom{={}}180\div 6\\
&=180\div (2\times 3)\\
&=180\div 2\div 3\\
&=90\div 3\\
&=30\end{aligned}}
\MethodTitle{第4小题：$560\div 70$}

\MethodBlock{口算思路一：同时去掉相同个数的0}{}{\begin{aligned}[t]&\phantom{={}}560\div 70\\
&=(560\div 10)\div (70\div 10)\\
&=56\div 7\\
&=8\end{aligned}}

\MethodBlock{口算思路二：拆分除数，连续相除}{}{\begin{aligned}[t]&\phantom{={}}560\div 70\\
&=560\div (2\times 35)\\
&=560\div 2\div 35\\
&=280\div 35\\
&=8\end{aligned}}
\MethodTitle{第5小题：$450\div 9$}

\MethodBlock{口算思路一：拆分除数，连续相除}{}{\begin{aligned}[t]&\phantom{={}}450\div 9\\
&=450\div (3\times 3)\\
&=450\div 3\div 3\\
&=150\div 3\\
&=50\end{aligned}}
\MethodTitle{第6小题：$15\times 6$}

\MethodBlock{口算思路一：按数位拆分计算}{把一个因数拆成整十、整百数与余下部分之和。}{\begin{aligned}[t]&\phantom{={}}15\times 6\\
&=6\times (10+5)\\
&=6\times 10+6\times 5\\
&=60+30\\
&=90\end{aligned}}
\MethodTitle{第7小题：$720\div 80$}

\MethodBlock{口算思路一：同时去掉相同个数的0}{}{\begin{aligned}[t]&\phantom{={}}720\div 80\\
&=(720\div 10)\div (80\div 10)\\
&=72\div 8\\
&=9\end{aligned}}

\MethodBlock{口算思路二：拆分除数，连续相除}{}{\begin{aligned}[t]&\phantom{={}}720\div 80\\
&=720\div (2\times 40)\\
&=720\div 2\div 40\\
&=360\div 40\\
&=9\end{aligned}}
\MethodTitle{第8小题：$20\times 50$}

\MethodBlock{口算思路一：拆因数，凑整计算}{把一个因数拆成两个因数，先凑成整百或整千。}{\begin{aligned}[t]&\phantom{={}}20\times 50\\
&=50\times (2\times 10)\\
&=(50\times 2)\times 10\\
&=100\times 10\\
&=1000\end{aligned}}

\MethodBlock{口算思路二：把特殊因数改成整百或整千}{先把25、125等改写成整百或整千除以一个数。}{\begin{aligned}[t]&\phantom{={}}20\times 50\\
&=(100\div 2)\times 20\\
&=100\times (20\div 2)\\
&=100\times 10\\
&=1000\end{aligned}}
\MethodTitle{第9小题：$360\div 40$}

\MethodBlock{口算思路一：同时去掉相同个数的0}{}{\begin{aligned}[t]&\phantom{={}}360\div 40\\
&=(360\div 10)\div (40\div 10)\\
&=36\div 4\\
&=9\end{aligned}}

\MethodBlock{口算思路二：拆分除数，连续相除}{}{\begin{aligned}[t]&\phantom{={}}360\div 40\\
&=360\div (2\times 20)\\
&=360\div 2\div 20\\
&=180\div 20\\
&=9\end{aligned}}
\MethodTitle{第10小题：$18\times 5$}

\MethodBlock{口算思路一：凑整后补偿}{把接近整十、整百的因数凑整，再补回或减去多算的部分。}{\begin{aligned}[t]&\phantom{={}}18\times 5\\
&=5\times (20-2)\\
&=5\times 20-5\times 2\\
&=100-10\\
&=90\end{aligned}}

\MethodBlock{口算思路二：按数位拆分计算}{把一个因数拆成整十、整百数与余下部分之和。}{\begin{aligned}[t]&\phantom{={}}18\times 5\\
&=5\times (10+8)\\
&=5\times 10+5\times 8\\
&=50+40\\
&=90\end{aligned}}
\MethodTitle{第11小题：$60\times 13$}

\MethodBlock{口算思路一：按数位拆分计算}{把一个因数拆成整十、整百数与余下部分之和。}{\begin{aligned}[t]&\phantom{={}}60\times 13\\
&=60\times (10+3)\\
&=60\times 10+60\times 3\\
&=600+180\\
&=780\end{aligned}}
\MethodTitle{第12小题：$480\div 6$}

\MethodBlock{口算思路一：利用商不变的性质}{被除数和除数同时乘同一个非零数，商不变。}{\begin{aligned}[t]&\phantom{={}}480\div 6\\
&=(480\times 5)\div (6\times 5)\\
&=2400\div 30\\
&=80\end{aligned}}

\MethodBlock{口算思路二：拆分除数，连续相除}{}{\begin{aligned}[t]&\phantom{={}}480\div 6\\
&=480\div (2\times 3)\\
&=480\div 2\div 3\\
&=240\div 3\\
&=80\end{aligned}}


\par\noindent

\QHeading{第2题　脱式计算。}

\CalcBlock[方法一]{1}{\begin{aligned}[t]&\phantom{={}}\color{black}(45+35)\times2\div16\\&\color{answerred}=80\times2\div16\\&\color{answerred}=160\div16\\&\color{answerred}=10\end{aligned}}
\par
\MethodBlock{方法二：先把能整除的两个数相除}{乘一个数再除以另一个数，可以先把能够整除的部分算出来；这里明确写出等价变形。}{\begin{aligned}[t]&\phantom{={}}(45+35)\times 2\div 16\\
&=80\times 2\div 16\\
&=[(80\div 16)\times 2]\\
&=[5\times 2]\\
&=10\end{aligned}}


\CalcBlock[方法一]{2}{\begin{aligned}[t]&\phantom{={}}\color{black}480\div[20\times(15-11)]\\&\color{answerred}=480\div[20\times4]\\&\color{answerred}=480\div80\\&\color{answerred}=6\end{aligned}}
\MethodBlock{方法二：除以积，改成连续除}{连续除以两个数，等于除以这两个数的积。}{\begin{aligned}[t]&\phantom{={}}480\div [20\times (15-11)]\\
&=480\div 20\div (15-11)\\
&=480\div 20\div 4\\
&=24\div 4\\
&=6\end{aligned}}

\CalcBlock{3}{\begin{aligned}[t]&\phantom{={}}\color{black}150\times[(84-39)\div9]\\&\color{answerred}=150\times[45\div9]\\&\color{answerred}=150\times5\\&\color{answerred}=750\end{aligned}}
\CalcBlock[方法一]{4}{\begin{aligned}[t]&\phantom{={}}\color{black}720\div[(12+6)\times4]\\&\color{answerred}=720\div[18\times4]\\&\color{answerred}=720\div72\\&\color{answerred}=10\end{aligned}}
\MethodBlock{方法二：除以积，改成连续除}{连续除以两个数，等于除以这两个数的积。}{\begin{aligned}[t]&\phantom{={}}720\div [(12+6)\times 4]\\
&=720\div (12+6)\div 4\\
&=720\div 18\div 4\\
&=40\div 4\\
&=10\end{aligned}}

\CalcBlock[方法一]{5}{\begin{aligned}[t]&\phantom{={}}\color{black}[280-(120+80)]\times5\\&\color{answerred}=[280-200]\times5\\&\color{answerred}=80\times5\\&\color{answerred}=400\end{aligned}}
\MethodBlock{方法二：乘法分配律}{分别求出两部分的积，再合起来。}{\begin{aligned}[t]&\phantom{={}}[280-(120+80)]\times 5\\
&=280\times 5-(120+80)\times 5\\
&=280\times 5-200\times 5\\
&=1400-1000\\
&=400\end{aligned}}

\CalcBlock[方法一]{6}{\begin{aligned}[t]&\phantom{={}}\color{black}960\div[(36-28)\times10]\\&\color{answerred}=960\div[8\times10]\\&\color{answerred}=960\div80\\&\color{answerred}=12\end{aligned}}
\MethodBlock{方法二：除以积，改成连续除}{连续除以两个数，等于除以这两个数的积。}{\begin{aligned}[t]&\phantom{={}}960\div [(36-28)\times 10]\\
&=960\div (36-28)\div 10\\
&=960\div 8\div 10\\
&=120\div 10\\
&=12\end{aligned}}

\par\noindent

\QHeading{第3题　解决问题。}

\Problem{①某粮店购进一批粮食，其中大米有600千克，绿豆有550千克。如果每25千克装一袋，大米比绿豆多装多少袋？}
\par


\MethodTitle{方法一}
\Answer{\Formula{\begin{aligned}[t]&\phantom{={}}(600-550)\div25\\&=50\div25\\&=2\text{（袋）}\end{aligned}}\par 答：大米比绿豆多装2袋。}

\MethodTitle{方法二：分别求装袋数，再求差}
\Answer{两种粮食每袋都是25千克，可以先分别求袋数。}
\Answer{\Formula{\begin{aligned}[t]&\phantom{={}}600\div 25-550\div 25\\
&=24-22\\
&=2\text{（袋）}\end{aligned}}}
\Answer{答：大米比绿豆多装2袋。}
\Problem{②某书店购进了560本图书，现有5个书架，每个书架有4层，每层摆放图书22本。将书架全部摆满后还剩多少本图书没有摆放？}

\MethodTitle{方法一：简便计算}
\Answer{\Formula{\begin{aligned}[t]&\phantom{={}}560-22\times 4\times 5\\&=560-22\times (4\times 5)\\&=560-22\times 20\\&=560-440\\&=120\text{（本）}\end{aligned}}\par 答：还剩120本图书没有摆放。}

\MethodTitle{方法二：先求总层数，再求已摆放本数}
\Answer{把全部书架的层数合起来，再按每层22本计算。}
\Answer{\Formula{\begin{aligned}[t]&\phantom{={}}560-22\times (4\times 5)\\
&=560-22\times 20\\
&=560-440\\
&=120\text{（本）}\end{aligned}}}
\Answer{答：还剩120本图书没有摆放。}
